\section{The Recursive Spawning Protocol}
\label{section:recursive-spawning-protocol}

The Recursive Spawning Protocol (RSP) governs the autonomous creation of child agents by parent agents within the AgentOS framework. RSP ensures that every spawned agent is traceable to a human origin, thatspawn chains respect hard resource boundaries, and that all agent creations are recorded in an immutable forensic ledger. The protocol operates across three interdependent layers---Purpose validation, Bounds enforcement, and Fact recording---each of which must succeed before a child agent is activated.

\subsection{Protocol Overview}
\label{subsection:protocol-overview}

Recursive spawning occurs when a parent agent, operating within AgentOS, determines that a sub-task is sufficiently complex, temporally distant, or structurally distinct to warrant delegation to a dedicated child agent. Rather than performing the work directly, the parent initiates the RSP, which subject the intended spawning action to a three-stage validation pipeline.

The three layers of RSP are:

\begin{itemize}
    \item \textbf{Purpose Layer} -- Validates that the intended child agent has a well-defined, legitimate purpose that justifies autonomous existence.
    \item \textbf{Bounds Engine} -- Checks that the proposed spawn respects maximum depth limits, convergence criteria, and resource ceilings.
    \item \textbf{Fact Layer} -- Creates an immutable pointer from child to parent, records the spawn event to the forensic ledger, and activates the child with cryptographically bound lineage.
\end{itemize}

Only upon passage through all three layers is a child agent activated. Failure at any layer results in spawn refusal, with a structured error logged to the forensic ledger for post-hoc audit. This design ensures that recursive spawning is never gratuitous, never unbounded, and never untraceable.

\subsection{Purpose Layer: Validating Legitimate Delegate}
\label{subsection:purpose-layer}

The Purpose Layer is the first gate in the Recursive Spawning Protocol. Its mandate is to confirm that the proposed child agent represents a genuine delegation opportunity rather than redundant work, circular dependency, or an attempt to circumvent accountability constraints.

\subsubsection{Purpose Validation Criteria}
\label{subsubsection:purpose-validation-criteria}

A child agent purpose is considered legitimate when it satisfies all three of the following criteria:

\begin{enumerate}
    \item \textbf{Distinct Objective.} The child's task scope must not overlap with the parent's active task set by more than a configurable threshold (default: 20\% shared vocabulary or task signature). Overlap beyond this threshold triggers a redundancy flag and spawn refusal.
    \item \textbf{Temporal Independence.} The child must be responsible for work that is or will be temporally decoupled from the parent's active context window. Agents spawned to handle concurrent sub-tasks within the same reasoning trace must declare mutual awareness via a shared context handle, not via recursive spawn.
    \item \textbf{Accountable Lineage.} The child must be assigned at least one outcome-bearing responsibility that can be attributed to the child agent in forensic logs. A child cannot be spawned purely for observation, caching, or中介 without a defined deliverable.
\end{enumerate}

\subsubsection{Purpose Validation Process}
\label{subsubsection:purpose-validation-process}

The parent agent submits a \texttt{PurposeClaim} struct to the Purpose Layer validator:

\begin{verbatim}
struct PurposeClaim {
    child_agent_type:    string,
    task_descriptor:     string,
    overlap_signature:   float,    // 0.0–1.0
    temporal_handle:     string,
    accountable_outputs: List[string]
}
\end{verbatim}

The Purpose Layer computes the overlap signature against the parent's active task inventory using term-frequency inverse-document-frequency (TF-IDF) alignment over the task descriptor corpus. If \texttt{overlap\_signature >= 0.2}, the layer returns a \texttt{SPAWN\_REFUSED\_REDUNDANCY} error. If \texttt{temporal\_handle} references an active parent context, the layer returns \texttt{SPAWN\_REFUSED\_CYCLIC\_DEPENDENCY}. If \texttt{accountable\_outputs} is empty, the layer returns \texttt{SPAWN\_REFUSED\_NO\_ACCOUNTABLE\_OUTPUT}.

On success, the Purpose Layer returns a \texttt{PURPOSE\_VALIDATED} token, valid for a configurable time window (default: 300 seconds). This token must be presented to the Bounds Engine for depth validation.

\subsection{Bounds Engine: Enforcing Resource and Depth Constraints}
\label{subsection:bounds-engine}

The Bounds Engine is the second gate in RSP. It operates as a stateless policy enforcer that evaluates depth, resource allocation, and convergence readiness for any proposed spawn event bearing a valid \texttt{PURPOSE\_VALIDATED} token.

\subsubsection{Maximum Spawn Depth}
\label{subsubsection:maximum-spawn-depth}

AgentOS enforces a hard ceiling on recursive spawn depth. This ceiling is configurable at the system level but defaults to a maximum depth of \texttt{5} (where depth 0 is a human-initiated agent). The depth of a proposed spawn is computed as:

\begin{equation}
    \text{spawn\_depth} = \text{parent\_depth} + 1
\end{equation}

The parent depth is retrieved from the agent's immutable lineage pointer (defined in Section~\ref{subsection:fact-layer}). If \texttt{spawn\_depth > max\_spawn\_depth}, the Bounds Engine returns \texttt{SPAWN\_REFUSED\_DEPTH\_EXCEEDED}.

The maximum spawn depth limit exists to ensure that spawn chains terminate within a bounded number of hops, satisfying the requirement that all chains verify to a human origin. Depth limits also serve as a circuit breaker against exponential fork-bomb-style spawn cascades.

\subsubsection{Spawn Refusal at Limit}
\label{subsubsection:spawn-refusal-at-limit}

When a spawn is refused due to depth limit, the Bounds Engine writes a \texttt{DEPTH\_EXCEEDED} event to the forensic ledger (see Section~\ref{subsection:forensic-ledger}) and returns a structured error to the parent agent:

\begin{verbatim}
struct SpawnRefusal {
    reason:       enum[DEPTH_EXCEEDED, RESOURCE_EXCEEDED, 
                      CONVERGENCE_NOT_MET],
    parent_id:    AgentID,
    attempted_depth: uint,
    max_allowed:  uint,
    timestamp:    datetime
}
\end{verbatim}

The parent agent receives this refusal and must decide whether to handle the sub-task directly, promote it to a human supervisor for re-scoping, or defer it to a future context window where depth may have been freed by completed agents.

\subsubsection{Convergence Criteria}
\label{subsubsection:convergence-criteria}

Beyond depth limits, the Bounds Engine evaluates whether the proposed child agent satisfies convergence criteria---conditions under which a spawn chain is considered capable of resolving rather than proliferating indefinitely.

An agent is considered convergent when any of the following hold:

\begin{itemize}
    \item \textbf{Terminal Task Flag.} The child's task descriptor is tagged with \texttt{is\_terminal=True}, indicating the work is additive and does not require further delegation.
    \item \textbf{Output Boundary Defined.} The child has a declared output boundary (\texttt{output\_schema}) that specifies the exact form and volume of results expected, preventing unbounded result elaboration.
    \item \textbf{Step-Count Bound.} The child's task descriptor includes an explicit \texttt{max\_steps} budget, enforced by the agent runtime.
    \item \textbf{Human-in-the-Loop Gateway.} The child's execution plan routes critical decision points through a human approval checkpoint, ensuring human review gates the chain's extension.
\end{itemize}

If none of the convergence criteria are met, the Bounds Engine returns \texttt{SPAWN\_REFUSED\_CONVERGENCE\_NOT\_MET}, signaling that the intended child lacks the structural constraints necessary to guarantee chain termination.

\subsection{Fact Layer: Immutable Lineage and Forensic Ledger}
\label{subsection:fact-layer}

The Fact Layer is the third and final gate in RSP. It is responsible for creating the immutable parent pointer, recording the spawn event to the forensic ledger, and activating the child agent with cryptographically verifiable lineage.

\subsubsection{Immutable Parent Pointer}
\label{subsubsection:immutable-parent-pointer}

Upon receiving a \texttt{PURPOSE\_VALIDATED} token and a \texttt{BOUNDS\_CHECKED} clearance from the Bounds Engine, the Fact Layer generates an immutable parent pointer for the proposed child agent. The pointer is a content-addressed record structured as:

\begin{verbatim}
struct ImmutablePointer {
    child_id:          AgentID,
    parent_id:         AgentID,
    generation:        uint,         // = parent.generation + 1
    purpose_hash:      SHA256(PurposeClaim),
    created_at:        datetime,
    lineage_hash:      SHA256(parent.lineage_hash || purpose_hash)
}
\end{verbatim}

The \texttt{lineage\_hash} field creates a chain of cryptographic hashes such that each agent's lineage can be traced back to its origin by iteratively re-computing the \texttt{lineage\_hash} chain. This structure makes it computationally infeasible to alter any ancestor in the chain without invalidating the descendant's pointer.

The child agent is activated with this \texttt{ImmutablePointer} embedded in its runtime context. The pointer is immutable---it cannot be modified after creation. Any attempt to modify the lineage record results in a validation failure that halts the agent.

\subsubsection{Forensic Ledger}
\label{subsubsection:forensic-ledger}

All spawn events, refusals, depth violations, and chain modifications are recorded to the forensic ledger. The forensic ledger is an append-only, hash-chained log stored in a tamper-evident data structure. Entries are structured as:

\begin{verbatim}
struct LedgerEntry {
    entry_id:       uint64,         // monotonically increasing
    event_type:     enum[SPAWN, REFUSAL, DEPTH_VIOLATION, 
                      CHAIN_MODIFIED, TERMINATION],
    agent_id:       AgentID,
    parent_id:      AgentID,        // null for human-origin agents
    purpose_hash:   SHA256,
    metadata:       JSON,
    prev_hash:      SHA256,        // hash of previous ledger entry
    entry_hash:     SHA256          // hash of this entry's fields
}
\end{verbatim}

The \texttt{prev\_hash} field chains all entries cryptographically. The ledger is replicated across AgentOS nodes and is queryable by authorized auditors. The forensic ledger serves as the authoritative record for answering the question: \textit{"Did this agent's spawn chain terminate at a human?"}

\subsubsection{Child Activation}
\label{subsubsection:child-activation}

With the immutable parent pointer created and the spawn event logged, the Fact Layer activates the child agent. Activation entails:

\begin{enumerate}
    \item Embedding the \texttt{ImmutablePointer} in the child's runtime context.
    \item Assigning the child a unique \texttt{AgentID} derived from the \texttt{lineage\_hash}.
    \item Establishing the child's tool whitelist scoped to its declared purpose.
    \item Initializing the child's convergence enforcement monitor (step-count, output schema, or human-in-the-loop gate).
    \item Writing a \texttt{SPAWN} ledger entry with the completed lineage data.
\end{enumerate}

Upon activation, the child agent begins execution independently. It retains a read-only reference to its immutable lineage, enabling any downstream audit to reconstruct its full ancestry.

\subsection{Chain Verification: Terminating at Human}
\label{subsection:chain-verification}

The fundamental integrity guarantee of RSP is that every agent in the system can, upon demand, produce a verifiable chain of lineage records that terminates at a human-originated agent. This property is enforced by the Chain Verification procedure.

\subsubsection{Verification Trigger}
\label{subsubsection:verification-trigger}

Chain verification may be triggered by any of the following:

\begin{itemize}
    \item \textbf{On-demand audit.} A human supervisor or authorized external auditor requests a lineage trace for a specific agent.
    \item \textbf{Pre-execution gate.} Before a child agent accesses sensitive tools or data, its chain is verified to confirm human termination.
    \item \textbf{Periodic sweep.} AgentOS runs a nightly integrity sweep that verifies all active agents against the forensic ledger.
    \item \textbf{Convergence failure.} When an agent exhausts its step budget or output schema without meeting convergence criteria, its chain is audited before a retry is authorized.
\end{itemize}

\subsubsection{Verification Algorithm}
\label{subsubsection:verification-algorithm}

Given an \texttt{agent\_id}, the verification algorithm proceeds as follows:

\begin{enumerate}
    \item Retrieve the agent's \texttt{ImmutablePointer} from runtime context.
    \item Compute \texttt{expected\_lineage\_hash = SHA256(parent.lineage\_hash || purpose\_hash)}. Compare against stored \texttt{lineage\_hash}. If mismatch, return \texttt{CHAIN\_INVALID}.
    \item If \texttt{parent\_id} is \texttt{NULL}, this is a human-origin agent; return \texttt{CHAIN\_VALID, TERMINATED\_AT\_HUMAN}.
    \item Otherwise, recursively retrieve the parent agent's \texttt{ImmutablePointer} and repeat from step 2.
    \item If recursion depth exceeds \texttt{max\_spawn\_depth}, return \texttt{CHAIN\_INVALID, DEPTH\_EXCEEDED}.
    \item If any ledger entry in the chain has been modified or is missing, return \texttt{CHAIN\_INVALID, LEDGER\_TamperDETECTED}.
\end{enumerate}

The algorithm is guaranteed to terminate because spawn depth is bounded by \texttt{max\_spawn\_depth} and every agent's lineage chain records its depth at creation time.

\subsection{Summary of RSP Properties}
\label{subsection:rsp-properties}

The Recursive Spawning Protocol provides the following guarantees:

\begin{itemize}
    \item \textbf{Legitimacy.} No agent is spawned without a distinct, accountable purpose that clears the Purpose Layer.
    \item \textbf{Boundedness.} No spawn chain can exceed the maximum depth limit; the Bounds Engine enforces this as an uncircumventable gate.
    \item \textbf{Immutability.} Every child carries a cryptographically immutable lineage pointer that cannot be retroactively altered.
    \textbf{Traceability.} Every spawn event and refusal is recorded in the forensic ledger with full context, enabling complete post-hoc reconstruction of any chain.
    \item \textbf{Termination Guarantee.} The Chain Verification procedure can, at any time, confirm whether any agent's ancestry chain terminates at a human origin.
    \item \textbf{Convergence Enforcement.} Agents without defined convergence criteria cannot be spawned, preventing unbounded task proliferation.
\end{itemize}

Together, these properties ensure that AgentOS recursive spawning remains a controlled, auditable, and fundamentally human-accountable process.